\Chapter{7}

\Verse{1}
{
    \zA{话说周瑞家的送了刘老老去后，便上来回王夫人话，谁知王夫人不在上房，问 丫鬟们，方知往薛姨妈那边说话儿去了。}
}
{
    \eA{To resume our narrative. Chou Jui's wife having seen old goody Liu off,
        speedily came to report the visit to madame Wang; but, contrary to her
        expectation, she did not find madame Wang in the drawing-room; and it was
        after inquiring of the waiting-maids that she eventually learnt that she had
        just gone over to have a chat with "aunt" Hsüeh.}
}
{
}

\Verse{2}
{
    \zA{周瑞家的听说，便出东角门过东院往梨香 院来。}
}
{
    \eA{Mrs. Chou, upon hearing this, hastily went out by the eastern corner door,
        and through the yard on the east,}
}
{
}

\Verse{3}
{
    \zA{刚至院门前，只见王夫人的丫鬟金钏儿和那一个才留头的小女孩儿站在台阶 儿上玩呢。}
}
{
    \eA{into the Pear Fragrance Court. As soon as she reached the entrance, she
        caught sight of madame Wang's waiting-maid, Chin Ch'uan-erh, playing about
        on the terrace steps, with a young girl, who had just let her hair grow.}
}
{
}

\Verse{4}
{
    \zA{看见周瑞家的进来，便知有话来回，因往里努嘴儿。}
}
{
    \eA{When they saw Chou Jui's wife approach, they forthwith surmised that she
        must have some message to deliver,}
}
{
}

\Verse{5}
{
    \zA{周瑞家的轻轻掀帘进去，见王夫人正和薛姨妈长篇大套的说些家务人情话。}
}
{
    \eA{so they pursed up their lips and directed her to the inner-room. Chou Jui's
        wife gently raised the curtain-screen, and upon entering discovered madame
        Wang, in voluble conversation with "aunt" Hsüeh,}
}
{
}

\Verse{6}
{
    \zA{周 瑞家的不敢惊动，遂进里间来。}
}
{
    \eA{about family questions and people in general. Mrs. Chou did not venture to
        disturb them,}
}
{
}

\Verse{7}
{
    \zA{只见薛宝钗家常打扮，头上只挽着？}
}
{
    \eA{and accordingly came into the inner room, where she found Hsüeh Pao-ch'ai in
        a house dress,}
}
{
}

\Verse{8}
{
    \zA{儿，坐在炕里 边，伏在几上和丫鬟莺儿正在那里描花样子呢。}
}
{
    \eA{with her hair simply twisted into a knot round the top of the head, sitting
        on the inner edge of the stove-couch, leaning on a small divan table,}
}
{
}

\Verse{9}
{
    \zA{见他进来，便放下笔，转过身，满 面堆笑让：“周姐姐坐。”}
}
{
    \eA{in the act of copying a pattern for embroidery, with the waiting-maid Ying
        Erh. When she saw her enter, Pao Ch'ai hastily put down her pencil,}
}
{
}

\Verse{10}
{
    \zA{周瑞家的也忙陪笑问道：“姑娘好？”}
}
{
    \eA{and turning round with a face beaming with smiles, "Sister Chou," she said,}
}
{
}

\Verse{11}
{
    \zA{一面炕沿边坐了， 因说：“这有两三天也没见姑娘到那边逛逛去，只怕是你宝兄弟冲撞了你不成？”}
}
{
    \eA{"take a seat." Chou Jui's wife likewise promptly returned the smile. "How is
        my young lady?" she inquired, as she sat down on the edge of the couch. "I
        haven't seen you come over on the other side for two or three days!}
}
{
}

\Verse{12}
{
    \zA{宝钗笑道：“那里的话。}
}
{
    \eA{Has Mr. Pao-yü perhaps given you offence?" "What an idea!" exclaimed Pao
        Ch'ai,}
}
{
}

\Verse{13}
{
    \zA{只因我那宗病又发了，所以且静养两天。”}
}
{
    \eA{with a smile. "It's simply that I've had for the last couple of days my old
        complaint again, and that I've in consequence kept quiet all this time,}
}
{
}

\Verse{14}
{
    \zA{周瑞家的道： “正是呢。}
}
{
    \eA{and looked after myself." "Is that it?" asked Chou Jui's wife;}
}
{
}

\Verse{15}
{
    \zA{姑娘到底有什么病根儿？}
}
{
    \eA{"but after all, what rooted kind of complaint are you subject to,}
}
{
}

\Verse{16}
{
    \zA{也该趁早请个大夫认真医治医治。}
}
{
    \eA{miss? you should lose really no time in sending for a doctor to diagnose it,
        and give you something to make you all right.}
}
{
}

\Verse{17}
{
    \zA{小小的年纪 儿倒作下个病根儿，也不是玩的呢。”}
}
{
    \eA{With your tender years, to have an organic ailment is indeed no trifle!" Pao
        Ch'ai laughed when she heard these remarks.}
}
{
}

\Verse{18}
{
    \zA{宝钗听说笑道：“再别提起这个病!也不知 请了多少大夫，吃了多少药，花了多少钱，总不见一点效验儿。}
}
{
    \eA{"Pray," she said, "don't allude to this again; for this ailment of mine I've
        seen, I can't tell you, how many doctors; taken no end of medicine and spent
        I don't know how much money; but the more we did so, not the least little
        bit of relief did I see.}
}
{
}

\Verse{19}
{
    \zA{后来还亏了一个和 尚，专治无名的病症，因请他看了。}
}
{
    \eA{Lucky enough, we eventually came across a bald-pated bonze, whose speciality
        was the cure of nameless illnesses. We therefore sent for him to see me,}
}
{
}

\Verse{20}
{
    \zA{他说我这是从胎里带来的一股热毒，幸而我先 天壮还不相干，要是吃凡药是不中用的。}
}
{
    \eA{and he said that I had brought this along with me from the womb as a sort of
        inflammatory virus, that luckily I had a constitution strong and hale so
        that it didn't matter; and that it would be of no avail if I took pills or
        any medicines. He then told me a prescription from abroad,}
}
{
}

\Verse{21}
{
    \zA{他就说了个海上仙方儿，又给了一包末药 作引子，异香异气的。}
}
{
    \eA{and gave me also a packet of a certain powder as a preparative, with a
        peculiar smell and strange flavour. He advised me, whenever my complaint
        broke out, to take a pill, which would be sure to put me right again.}
}
{
}

\Verse{22}
{
    \zA{他说犯了时吃一丸就好了。}
}
{
    \eA{And this has, after all, strange to say, done me a great deal of good."}
}
{
}

\Verse{23}
{
    \zA{倒也奇怪，这倒效验些。”}
}
{
    \eA{"What kind of prescription is this one from abroad, I wonder," remarked Mrs.
        Chou;}
}
{
}

\Verse{24}
{
    \zA{周瑞 家的因问道：“不知是什么方儿？}
}
{
    \eA{"if you, miss, would only tell me, it would be worth our while bearing it in
        mind, and recommending it to others:}
}
{
}

\Verse{25}
{
    \zA{姑娘说了，我们也好记着说给人知道。}
}
{
    \eA{and if ever we came across any one afflicted with this disease, we would
        also be doing a charitable deed."}
}
{
}

\Verse{26}
{
    \zA{要遇见这 样病，也是行好的事”宝钗笑道：“不问这方儿还好，若问这方儿，真把人琐碎死
        了!东西药料一概却都有限，最难得是‘可巧’二字：要春天开的白牡丹花蕊十二 两，夏天开的白荷花蕊十二两，秋天的白芙蓉蕊十二两，冬天的白梅花蕊十二两。}
}
{
    \eA{"You'd better not ask for the prescription," rejoined Pao Ch'ai smiling.
        "Why, its enough to wear one out with perplexity! the necessaries and
        ingredients are few, and all easy to get, but it would be difficult to find
        the lucky moment! You want twelve ounces of the pollen of the white peone,
        which flowers in spring,}
}
{
}

\Verse{27}
{
    \zA{将这四样花蕊于次年春分这一天晒干，和在末药一处，一齐研好；又要雨水这日的 天落水十二钱……”周瑞家的笑道：“嗳呀，这么说就得三年的工夫呢。}
}
{
    \eA{twelve ounces of the pollen of the white summer lily, twelve ounces of the
        pollen of the autumn hibiscus flower, and twelve ounces of the white plum in
        bloom in the winter. You take the four kinds of pollen, and put them in the
        sun, on the very day of the vernal equinox of the succeeding year to get
        dry, and then you mix them with the powder and pound them well together.}
}
{
}

\Verse{28}
{
    \zA{倘或雨水 这日不下雨，可又怎么着呢？”}
}
{
    \eA{You again want twelve mace of water, fallen on 'rain water' day....." "Good
        gracious!" exclaimed Mrs. Chou promptly,}
}
{
}

\Verse{29}
{
    \zA{宝钗笑道：“所以了!那里有这么可巧的雨？}
}
{
    \eA{as she laughed. "From all you say, why you want three years' time! and what
        if no rain falls on 'rain water' day!}
}
{
}

\Verse{30}
{
    \zA{也只好 再等罢了。}
}
{
    \eA{What would one then do?" "Quite so!" Pao Ch'ai remarked smilingly;}
}
{
}

\Verse{31}
{
    \zA{还要白露这日的露水十二钱，霜降这日的霜十二钱，小雪这日的雪十二 钱。}
}
{
    \eA{"how can there be such an opportune rain on that very day! but to wait is
        also the best thing, there's nothing else to be done. Besides, you want
        twelve mace of dew, collected on 'White Dew' day,}
}
{
}

\Verse{32}
{
    \zA{把这四样水调匀了，丸了龙眼大的丸子，盛在旧磁坛里，埋在花根底下。}
}
{
    \eA{and twelve mace of the hoar frost, gathered on 'Frost Descent' day, and
        twelve mace of snow, fallen on 'Slight Snow' day! You next take these four
        kinds of waters and mix them with the other ingredients,}
}
{
}

\Verse{33}
{
    \zA{若发 了病的时候儿，拿出来吃一丸，用一钱二分黄柏煎汤送下。”}
}
{
    \eA{and make pills of the size of a lungngan. You keep them in an old porcelain
        jar, and bury them under the roots of some flowers; and when the ailment
        betrays itself,}
}
{
}

\Verse{34}
{
    \zA{周瑞家的听了，笑道：“阿弥陀佛!真巧死了人。}
}
{
    \eA{you produce it and take a pill, washing it down with two candareens of a
        yellow cedar decoction." "O-mi-to-fu!" cried Mrs. Chou,}
}
{
}

\Verse{35}
{
    \zA{等十年还未必碰的全呢！”}
}
{
    \eA{when she heard all this, bursting out laughing. "It's really enough to kill
        one!}
}
{
}

\Verse{36}
{
    \zA{宝钗道：“竟好。}
}
{
    \eA{you might wait ten years and find no such lucky moments!"}
}
{
}

\Verse{37}
{
    \zA{自他去后，一二年间，可巧都得了，好容易配成一料。}
}
{
    \eA{"Fortunate for me, however," pursued Pao Ch'ai, "in the course of a year or
        two, after the bonze had told me about this prescription,}
}
{
}

\Verse{38}
{
    \zA{如今从家 里带了来，现埋在梨花树底下。”}
}
{
    \eA{we got all the ingredients; and, after much trouble, we compounded a supply,
        which we have now brought along with us from the south to the north;}
}
{
}

\Verse{39}
{
    \zA{周瑞家的又道：“这药有名字没有呢？”}
}
{
    \eA{and lies at present under the pear trees." "Has this medicine any name or
        other of its own?" further inquired Mrs. Chou.}
}
{
}

\Verse{40}
{
    \zA{宝钗道： “有。}
}
{
    \eA{"It has a name," replied Pao Ch'ai;}
}
{
}

\Verse{41}
{
    \zA{也是那和尚说的，叫作‘冷香丸’。”}
}
{
    \eA{"the mangy-headed bonze also told it me; he called it 'cold fragrance'
        pill."}
}
{
}

\Verse{42}
{
    \zA{周瑞家的听了点头儿，因又说：“这 病发了时，到底怎么着？”}
}
{
    \eA{Chou Jui's wife nodded her head, as she heard these words. "What do you feel
        like after all when this complaint manifests itself?"}
}
{
}

\Verse{43}
{
    \zA{宝钗道：“也不觉什么，不过只喘嗽些，吃一丸也就罢 了。”}
}
{
    \eA{she went on to ask. "Nothing much," replied Pao Ch'ai; "I simply pant and
        cough a bit; but after I've taken a pill,}
}
{
}

\Verse{44}
{
    \zA{周瑞家的还要说话时，忽听王夫人问道：“谁在里头？”}
}
{
    \eA{I get over it, and it's all gone." Mrs. Chou was bent upon making some
        further remark, when madame Wang was suddenly heard to enquire,}
}
{
}

\Verse{45}
{
    \zA{周瑞家的忙出来答应 了，便回了刘老老之事。}
}
{
    \eA{"Who is in here?" Mrs. Chou went out hurriedly and answered; and forthwith
        told her all about old goody Liu's visit.}
}
{
}

\Verse{46}
{
    \zA{略待半刻，见王夫人无话，方欲退出去，薛姨妈忽又笑道： “你且站住。}
}
{
    \eA{Having waited for a while, and seeing that madame Wang had nothing to say,
        she was on the point of retiring, when "aunt" Hsueh unexpectedly remarked
        smiling:}
}
{
}

\Verse{47}
{
    \zA{我有一件东西，你带了去罢。”}
}
{
    \eA{"Wait a bit! I've something to give you to take along with you."}
}
{
}

\Verse{48}
{
    \zA{说着便叫：“香菱！”}
}
{
    \eA{And as she spoke, she called for Hsiang Ling.}
}
{
}

\Verse{49}
{
    \zA{帘栊响处，才 和金钏儿玩的那个小丫头进来，问：“太太叫我做什么？”}
}
{
    \eA{The sound of the screen-board against the sides of the door was heard, and
        in walked the waiting-maid, who had been playing with Chin Ch'uan-erh.}
}
{
}

\Verse{50}
{
    \zA{薛姨妈道：“把那匣子 里的花儿拿来。”}
}
{
    \eA{"Did my lady call?" she asked. "Bring that box of flowers," said Mrs. Hsueh.
        Hsiang Ling assented,}
}
{
}

\Verse{51}
{
    \zA{香菱答应了，向那边捧了个小锦匣儿来。}
}
{
    \eA{and brought from the other side a small embroidered silk box. "These,"
        explained "aunt" Hsüeh,}
}
{
}

\Verse{52}
{
    \zA{薛姨妈道：“这是宫里 头作的新鲜花样儿堆纱花，十二枝。}
}
{
    \eA{"are a new kind of flowers, made in the palace. They consist of twelve twigs
        of flowers of piled gauze. I thought of them yesterday,}
}
{
}

\Verse{53}
{
    \zA{昨儿我想起来，白放着可惜旧了，何不给他们 姐妹们戴去。}
}
{
    \eA{and as they will, the pity is, only get old, if uselessly put away, why not
        give them to the girls to wear them in their hair! I meant to have sent them
        over yesterday,}
}
{
}
\ChapterImage{img/chapters/021第7回 寶釵談病配冷香丸 薛姨嫣托周瑞送花 金釧周瑞笑香菱形.jpg}


\Verse{54}
{
    \zA{昨儿要送去，偏又忘了；你今儿来得巧，就带了去罢。}
}
{
    \eA{but I forgot all about them. You come to-day most opportunely, and if you
        will take them with you, I shall have got them off my hands.}
}
{
}

\Verse{55}
{
    \zA{你家的三位姑 娘每位两枝，下剩六枝送林姑娘两枝，那四枝给凤姐儿罢。”}
}
{
    \eA{To the three young ladies in your family give two twigs each, and of the six
        that will remain give a couple to Miss Lin, and the other four to lady
        Feng." "Better keep them and give them to your daughter Pao Ch'ai to wear,"}
}
{
}

\Verse{56}
{
    \zA{王夫人道：“留着给 宝丫头戴也罢了，又想着他们。”}
}
{
    \eA{observed madame Wang, "and have done with it; why think of all the others?"
        "You don't know, sister,"}
}
{
}

\Verse{57}
{
    \zA{薛姨妈道：“姨太太不知，宝丫头怪着呢，他从 来不爱这些花儿粉儿的。”}
}
{
    \eA{replied "aunt" Hsüeh, "what a crotchety thing Pao Ch'ai is! she has no
        liking for flower or powder." With these words on her lips, Chou Jui's wife
        took the box and walked out of the door of the room.}
}
{
}

\Verse{58}
{
    \zA{说着，周瑞家的拿了匣子，走出房门。}
}
{
    \eA{Perceiving that Chin Ch'uan-erh was still sunning herself outside, Chou
        Jui's wife asked her:}
}
{
}

\Verse{59}
{
    \zA{见金钏儿仍在那里晒日阳儿，周瑞家的 问道：“那香菱小丫头子可就是时常说的，临上京时买的、为他打人命官司的那个 小丫头吗？”}
}
{
    \eA{"Isn't this Hsiang Ling, the waiting-maid that we've often heard of as
        having been purchased just before the departure of the Hsüeh family for the
        capital, and on whose account there occurred some case of manslaughter or
        other?" "Of course it's she," replied Chin Ch'uan.}
}
{
}

\Verse{60}
{
    \zA{金钏儿道：“可不就是他。”}
}
{
    \eA{But as they were talking, they saw Hsiang Ling draw near smirkingly,}
}
{
}

\Verse{61}
{
    \zA{正说着，只见香菱笑嘻嘻的走来，周瑞 家的便拉了他的手细细的看了一回，因向金钏儿笑道：“这个模样儿，竟有些像咱 们东府里的小蓉奶奶的品格儿。”}
}
{
    \eA{and Chou Jui's wife at once seized her by the hand, and after minutely
        scrutinizing her face for a time, she turned round to Chin Ch'uan-erh and
        smiled. "With these features she really resembles slightly the style of lady
        Jung of our Eastern Mansion." "So I too maintain!" said Chin Ch'uan-erh.
        Chou Jui's wife then asked Hsiang Ling,}
}
{
}

\Verse{62}
{
    \zA{金钏儿道：“我也这么说呢。”}
}
{
    \eA{"At what age did you enter this family? and where are your father and mother
        at present?"}
}
{
}

\Verse{63}
{
    \zA{周瑞家的又问香 菱：“你几岁投身到这里？”}
}
{
    \eA{and also inquired, "In what year of your teens are you? and of what place
        are you a native?" But Hsiang Ling, after listening to all these questions,}
}
{
}

\Verse{64}
{
    \zA{又问：“你父母在那里呢？}
}
{
    \eA{simply nodded her head and replied, "I can't remember." When Mrs. Chou and
        Chin Ch'uan-erh heard these words,}
}
{
}

\Verse{65}
{
    \zA{今年十几了？}
}
{
    \eA{their spirits changed to grief,}
}
{
}

\Verse{66}
{
    \zA{本处是那里的 人？”}
}
{
    \eA{and for a while they felt affected and wounded at heart;}
}
{
}

\Verse{67}
{
    \zA{香菱听问，摇头说：“不记得了。”}
}
{
    \eA{but in a short time, Mrs. Chou brought the flowers into the room at the back
        of madame Wang's principal apartment. The fact is that dowager lady Chia had
        explained that as her granddaughters were too numerous, it would not be
        convenient to crowd them together in one place,}
}
{
}

\Verse{68}
{
    \zA{周瑞家的和金钏儿听了，倒反为叹息了 一回。}
}
{
    \eA{that Pao-yü and Tai-yü should only remain with her in this part to break her
        loneliness, but that Ying Ch'un, T'an Ch'un,}
}
{
}

\Verse{69}
{
    \zA{一时周瑞家的携花至王夫人正房后。}
}
{
    \eA{and Hsi Ch'un, the three of them, should move on this side in the three
        rooms within the antechamber,}
}
{
}

\Verse{70}
{
    \zA{原来近日贾母说孙女们太多，一处挤着倒 不便，只留宝玉黛玉二人在这边解闷，却将迎春、探春、惜春三人移到王夫人这边 房后三间抱厦内居住，令李纨陪伴照管。}
}
{
    \eA{at the back of madame lady Wang's quarters; and that Li Wan should be told
        off to be their attendant and to keep an eye over them. Chou Jui's wife,
        therefore, on this occasion came first to these rooms as they were on her
        way, but she only found a few waiting-maids assembled in the antechamber,
        waiting silently to obey a call.}
}
{
}

\Verse{71}
{
    \zA{如今周瑞家的故顺路先往这里来，只见几 个小丫头都在抱厦内默坐，听着呼唤。}
}
{
    \eA{Ying Ch'un's waiting-maid, Ssu Chi, together with Shih Shu, T'an Ch'un's
        waiting-maid, just at this moment raised the curtain, and made their egress,
        each holding in her hand a tea-cup and saucer;}
}
{
}

\Verse{72}
{
    \zA{迎春的丫鬟司棋和探春的丫鬟侍书二人，正 掀帘子出来，手里都捧着茶盘茶钟，周瑞家的便知他姐妹在一处坐着，也进入房内。}
}
{
    \eA{and Chou Jui's wife readily concluding that the young ladies were sitting
        together also walked into the inner room, where she only saw Ying Ch'un and
        T'an Ch'un seated near the window, in the act of playing chess. Mrs. Chou
        presented the flowers and explained whence they came,}
}
{
}

\Verse{73}
{
    \zA{只见迎春、探春二人正在窗下围棋。}
}
{
    \eA{and what they were. The girls forthwith interrupted their game, and both
        with a curtsey,}
}
{
}

\Verse{74}
{
    \zA{周瑞家的将花送上，说明原故，二人忙住了棋， 都欠身道谢，命丫鬟们收了。}
}
{
    \eA{expressed their thanks, and directed the waiting-maids to put the flowers
        away. Mrs. Chou complied with their wishes (and handing over the flowers);
        "Miss Hsi Ch'un," she remarked,}
}
{
}

\Verse{75}
{
    \zA{周瑞家的答应了，因说：“四姑娘不在房里，只怕在老太太那边呢？”}
}
{
    \eA{"is not at home; and possibly she's over there with our old lady." "She's in
        that room, isn't she?" inquired the waiting-maids. Mrs. Chou at these words
        readily came into the room on this side,}
}
{
}

\Verse{76}
{
    \zA{丫鬟们 道：“在那屋里不是？”}
}
{
    \eA{where she found Hsi Ch'un, in company with a certain Chih Neng,}
}
{
}

\Verse{77}
{
    \zA{周瑞家的听了，便往这边屋里来。}
}
{
    \eA{a young nun of the "moon reflected on water" convent, talking and laughing
        together.}
}
{
}

\Verse{78}
{
    \zA{只见惜春正同水月庵的 小姑子智能儿两个一处玩耍呢，见周瑞家的进来，便问他何事。}
}
{
    \eA{On seeing Chou Jui's wife enter, Hsi Ch'un at once asked what she wanted,
        whereupon Chou Jui's wife opened the box of flowers, and explained who had
        sent them. "I was just telling Chih Neng,"}
}
{
}

\Verse{79}
{
    \zA{周瑞家的将花匣打 开，说明原故，惜春笑道：“我这里正和智能儿说，我明儿也要剃了头跟他作姑子 去呢。}
}
{
    \eA{remarked Hsi Ch'un laughing, "that I also purpose shortly shaving my head
        and becoming a nun; and strange enough, here you again bring me flowers; but
        supposing I shave my head, where can I wear them?" They were all very much
        amused for a time with this remark,}
}
{
}

\Verse{80}
{
    \zA{可巧又送了花来，要剃了头，可把花儿戴在那里呢？”}
}
{
    \eA{and Hsi Ch'un told her waiting-maid, Ju Hua, to come and take over the
        flowers. "What time did you come over?" then inquired Mrs. Chou of Chih
        Neng.}
}
{
}

\Verse{81}
{
    \zA{说着，大家取笑一回， 惜春命丫鬟收了。}
}
{
    \eA{"Where is that bald-pated and crotchety superior of yours gone?" "We came,"
        explained Chih Neng,}
}
{
}

\Verse{82}
{
    \zA{周瑞家的因问智能儿：“你是什么时候来的？}
}
{
    \eA{"as soon as it was day; after calling upon madame Wang, my superior went
        over to pay a visit in the mansion of Mr. Yü,}
}
{
}

\Verse{83}
{
    \zA{你师父那秃歪剌那 里去了？”}
}
{
    \eA{and told me to wait for her here." "Have you received," further asked Mrs.
        Chou,}
}
{
}

\Verse{84}
{
    \zA{智能儿道：“我们一早就来了。}
}
{
    \eA{"the monthly allowance for incense offering due on the fifteenth or not?"}
}
{
}

\Verse{85}
{
    \zA{我师父见过太太，就往于老爷府里去了， 叫我在这里等他呢。”}
}
{
    \eA{"I can't say," replied Chih Neng. "Who's now in charge of the issue of the
        monthly allowances to the various temples?" interposed Hsi Ch'un, addressing
        Mrs. Chou,}
}
{
}

\Verse{86}
{
    \zA{周瑞家的又道：“十五的月例香供银子可得了没有？”}
}
{
    \eA{as soon as she heard what was said. "It's Yü Hsin," replied Chou Jui's wife,
        "who's intrusted with the charge."}
}
{
}

\Verse{87}
{
    \zA{智能 儿道：“不知道。”}
}
{
    \eA{"That's how it is," observed Hsi Ch'un with a chuckle;}
}
{
}

\Verse{88}
{
    \zA{惜春便问周瑞家的：“如今各庙月例银子是谁管着？”}
}
{
    \eA{"soon after the arrival of the Superior, Yü Hsin's wife came over and kept
        on whispering with her for some time; so I presume it must have been about
        this allowance."}
}
{
}

\Verse{89}
{
    \zA{周瑞家 的道：“余信管着。”}
}
{
    \eA{Mrs. Chou then went on to bandy a few words with Chih Neng, after which she
        came over to lady Feng's apartments.}
}
{
}

\Verse{90}
{
    \zA{惜春听了笑道：“这就是了。}
}
{
    \eA{Proceeding by a narrow passage, she passed under Li Wan's back windows,}
}
{
}

\Verse{91}
{
    \zA{他师父一来了，余信家的就赶 上来，和他师父咕唧了半日，想必就是为这个事了。”}
}
{
    \eA{and went along the wall ornamented with creepers on the west. Going out of
        the western side gate, she entered lady Feng's court, and walked over into
        the Entrance Hall, where she only found the waiting-girl Feng Erh,}
}
{
}

\Verse{92}
{
    \zA{那周瑞家的又和智能儿唠叨了一回，便往凤姐处来。}
}
{
    \eA{sitting on the doorsteps of lady Feng's apartments. When she caught sight of
        Mrs. Chou approaching, she at once waved her hand,}
}
{
}

\Verse{93}
{
    \zA{穿过了夹道子，从李纨后 窗下越过西花墙，出西角门，进凤姐院中。}
}
{
    \eA{bidding her go to the eastern room. Chou Jui's wife understood her meaning,
        and hastily came on tiptoe to the chamber on the east, where she saw a nurse
        patting lady Feng's daughter to sleep.}
}
{
}

\Verse{94}
{
    \zA{走至堂屋，只见小丫头丰儿坐在房门槛 儿上，见周瑞家的来了，连忙的摆手儿，叫他往东屋里去。}
}
{
    \eA{Mrs. Chou promptly asked the nurse in a low tone of voice: "Is the young
        lady asleep at this early hour? But if even she is I must wake her up." The
        nurse nodded her head in assent, but as these inquiries were being made,}
}
{
}

\Verse{95}
{
    \zA{周瑞家的会意，忙着蹑 手蹑脚儿的往东边屋里来，只见奶子拍着大姐儿睡觉呢。}
}
{
    \eA{a sound of laughter came from over the other side, in which lady Feng's
        voice could be detected; followed, shortly after, by the sound of a door
        opening, and out came P'ing Erh, with a large brass basin in her hands,}
}
{
}

\Verse{96}
{
    \zA{周瑞家的悄悄儿问道：“二 奶奶睡中觉呢吗？}
}
{
    \eA{which she told Feng Erh to fill with water and take inside. P'ing Erh
        forthwith entered the room on this side,}
}
{
}

\Verse{97}
{
    \zA{也该清醒了。”}
}
{
    \eA{and upon perceiving Chou Jui's wife:}
}
{
}

\Verse{98}
{
    \zA{奶子笑着，撇着嘴摇头儿。}
}
{
    \eA{"What have you come here again for, my old lady?" she readily inquired.}
}
{
}

\Verse{99}
{
    \zA{正问着，只听那边微 有笑声儿，却是贾琏的声音。}
}
{
    \eA{Chou Jui's wife rose without any delay, and handed her the box. "I've come,"
        said she, "to bring you a present of flowers."}
}
{
}

\Verse{100}
{
    \zA{接着房门响，平儿拿着大铜盆出来，叫人舀水。}
}
{
    \eA{Upon hearing this, P'ing Erh opened the box, and took out four sprigs, and,
        turning round, walked out of the room.}
}
{
}

\Verse{101}
{
    \zA{平儿 便进这边来，见了周瑞家的，便问：“你老人家又来作什么？”}
}
{
    \eA{In a short while she came from the inner room with two sprigs in her hand,
        and calling first of all Ts'ai Ming, she bade her take the flowers over to
        the mansion on the other side and present them to "madame" Jung,}
}
{
}

\Verse{102}
{
    \zA{周瑞家的忙起身拿 匣子给他看道：“送花儿来了。”}
}
{
    \eA{after which she asked Mrs. Chou to express her thanks on her return. Chou
        Jui's wife thereupon came over to dowager lady Chia's room on this side of
        the compound,}
}
{
}

\Verse{103}
{
    \zA{平儿听了，便打开匣子，拿了四枝，抽身去了。}
}
{
    \eA{and as she was going through the Entrance Hall, she casually came, face to
        face, with her daughter, got up in gala dress,}
}
{
}

\Verse{104}
{
    \zA{半刻工夫，手里拿出两枝来，先叫彩明来，吩咐：“送到那边府里，给小蓉大奶奶 戴的。”}
}
{
    \eA{just coming from the house of her mother-in-law. "What are you running over
        here for at this time?" promptly inquired Mrs. Chou. "Have you been well of
        late,}
}
{
}

\Verse{105}
{
    \zA{次后方命周瑞家的回去道谢。}
}
{
    \eA{mother?" asked her daughter. "I've been waiting for ever so long at home,}
}
{
}

\Verse{106}
{
    \zA{周瑞家的这才往贾母这边来，过了穿堂，顶头忽见他的女孩儿打扮着才从他婆 家来。}
}
{
    \eA{but you never come out! What's there so pressing that has prevented you from
        returning home? I waited till I was tired, and then went on all alone, and
        paid my respects to our venerable lady; I'm now, on my way to inquire about
        our lady Wang.}
}
{
}

\Verse{107}
{
    \zA{周瑞家的忙问：“你这会子跑来作什么？”}
}
{
    \eA{What errand haven't you delivered as yet, ma; and what is it you're
        holding?" "Ai! as luck would have it,"}
}
{
}
\ChapterImage{img/chapters/022第7回 周瑞送各姐妹宮花.jpg}


\Verse{108}
{
    \zA{他女孩儿说：“妈，一向身上好？}
}
{
    \eA{rejoined Chou Jui's wife smilingly, "old goody Liu came over to-day,}
}
{
}

\Verse{109}
{
    \zA{我在家里等了这半日，妈竟不去，什么事情这么忙的不回家？}
}
{
    \eA{so that besides my own hundred and one duties, I've had to run about here
        and there ever so long, and all for her! While attending to these,}
}
{
}

\Verse{110}
{
    \zA{我等烦了，自己先到 了老太太跟前请了安了，这会子请太太的安去。}
}
{
    \eA{Mrs. Hsueh came across me, and asked me to take these flowers to the young
        ladies, and I've been at it up to this very moment, and haven't done yet!
        But coming at this time,}
}
{
}

\Verse{111}
{
    \zA{妈还有什么不了的差事？}
}
{
    \eA{you must surely have something or other that you want me to do for you!}
}
{
}

\Verse{112}
{
    \zA{手里是什 么东西？”}
}
{
    \eA{what's it?" "Really ma, you're quick at guessing!"}
}
{
}

\Verse{113}
{
    \zA{周瑞家的笑道：“嗳!今儿偏偏来了个刘老老，我自己多事，为他跑了 半日。}
}
{
    \eA{exclaimed her daughter with a smile; "I'll tell you what it's all about. The
        day before yesterday, your son-in-law had a glass of wine too many, and
        began altercating with some person or other;}
}
{
}

\Verse{114}
{
    \zA{这会子叫姨太太看见了，叫送这几枝花儿给姑娘奶奶们去，这还没有送完呢。}
}
{
    \eA{and some one, I don't know why, spread some evil report, saying that his
        antecedents were not clear, and lodged a charge against him at the Yamen,
        pressing the authorities to deport him to his native place.}
}
{
}

\Verse{115}
{
    \zA{你今儿来，一定有什么事情。”}
}
{
    \eA{That's why I've come over to consult with you, as to whom we should appeal
        to,}
}
{
}

\Verse{116}
{
    \zA{他女孩儿笑道：“你老人家倒会猜，一猜就猜着了。}
}
{
    \eA{to do us this favour of helping us out of our dilemma!" "I knew at once,"
        Mrs. Chou remarked after listening, "that there was something wrong;}
}
{
}

\Verse{117}
{
    \zA{实对你老人家说：你女婿因前儿多喝了点子酒，和人分争起来，不知怎么叫人放了 把邪火，说他来历不明，告到衙门里，要递解还乡。}
}
{
    \eA{but this is nothing hard to settle! Go home and wait for me and I'll come
        straightway, as soon as I've taken these flowers to Miss Lin; our madame
        Wang and lady Secunda have both no leisure (to attend to you now,) so go
        back and wait for me!}
}
{
}

\Verse{118}
{
    \zA{所以我来和你老人家商量商量， 讨个情分。}
}
{
    \eA{What's the use of so much hurry!" Her daughter, upon hearing this, forthwith
        turned round to go back,}
}
{
}

\Verse{119}
{
    \zA{不知求那个可以了事？”}
}
{
    \eA{when she added as she walked away, "Mind, mother, and make haste."}
}
{
}

\Verse{120}
{
    \zA{周瑞家的听了道：“我就知道。}
}
{
    \eA{"All right," replied Chou Jui's wife, "of course I will; you are young yet,}
}
{
}

\Verse{121}
{
    \zA{这算什么大事， 忙的这么着!你先家去，等我送下林姑娘的花儿就回去。}
}
{
    \eA{and without experience, and that's why you are in this flurry." As she
        spoke, she betook herself into Tai-yü's apartments. Contrary to her
        expectation Tai-yü was not at this time in her own room,}
}
{
}

\Verse{122}
{
    \zA{这会儿太太二奶奶都不得 闲儿呢！”}
}
{
    \eA{but in Pao-yü's; where they were amusing themselves in trying to solve the
        "nine strung rings" puzzle. On entering Mrs. Chou put on a smile.}
}
{
}

\Verse{123}
{
    \zA{他女孩儿听说，便回去了，还说：“妈，好歹快来。”}
}
{
    \eA{"'Aunt' Hsüeh," she explained, "has told me to bring these flowers and
        present them to you to wear in your hair." "What flowers?" exclaimed Pao-yü.}
}
{
}

\Verse{124}
{
    \zA{周瑞家的道：“是 了罢!小人儿家没经过什么事，就急的这么个样儿。”}
}
{
    \eA{"Bring them here and let me see them." As he uttered these words, he readily
        stretched out his hands and took them over, and upon opening the box and
        looking in,}
}
{
}

\Verse{125}
{
    \zA{说着，便到黛玉房中去了。}
}
{
    \eA{he discovered, in fact, two twigs of a novel and artistic kind of artificial
        flowers,}
}
{
}

\Verse{126}
{
    \zA{谁知此时黛玉不在自己房里，却在宝玉房中，大家解九连环作戏。}
}
{
    \eA{of piled gauze, made in the palace. Tai-yü merely cast a glance at them, as
        Pao-yü held them. "Have these flowers," she inquired eagerly, "been sent to
        me alone,}
}
{
}

\Verse{127}
{
    \zA{周瑞家的进 来，笑道：“林姑娘，姨太太叫我送花儿来了。”}
}
{
    \eA{or have all the other girls got some too?" "Each one of the young ladies has
        the same," replied Mrs. Chou; "and these two twigs are intended for you,}
}
{
}

\Verse{128}
{
    \zA{宝玉听说，便说：“什么花儿？}
}
{
    \eA{miss." Tai-yü forced a smile. "Oh! I see," she observed.}
}
{
}

\Verse{129}
{
    \zA{拿来我瞧瞧。”}
}
{
    \eA{"If all the others hadn't chosen,}
}
{
}

\Verse{130}
{
    \zA{一面便伸手接过匣子来看时，原来是两枝宫制堆纱新巧的假花。}
}
{
    \eA{even these which remain over wouldn't have been given to me." Chou Jui's
        wife did not utter a word in reply. "Sister Chou, what took you over on the
        other side?"}
}
{
}

\Verse{131}
{
    \zA{黛 玉只就宝玉手中看了一看，便问道：“还是单送我一个人的，还是别的姑娘们都有 呢？”}
}
{
    \eA{asked Pao-yü. "I was told that our madame Wang was over there," explained
        Mrs. Chou, "and as I went to give her a message, 'aunt' Hsüeh seized the
        opportunity to ask me to bring over these flowers."}
}
{
}

\Verse{132}
{
    \zA{周瑞家的道：“各位都有了，这两枝是姑娘的。”}
}
{
    \eA{"What was cousin Pao Ch'ai doing at home?" asked Pao-yü. "How is it she's
        not even been over for these few days?"}
}
{
}

\Verse{133}
{
    \zA{黛玉冷笑道：“我就知道 么!别人不挑剩下的也不给我呀。”}
}
{
    \eA{"She's not quite well," remarked Mrs. Chou. When Pao-yü heard this news,
        "Who'll go," he speedily ascertained of the waiting-maids,}
}
{
}

\Verse{134}
{
    \zA{周瑞家的听了，一声儿也不敢言语。}
}
{
    \eA{"and inquire after her? Tell her that cousin Lin and I have sent round to
        ask how our aunt and cousin are getting on!}
}
{
}

\Verse{135}
{
    \zA{宝玉问道： “周姐姐，你作什么到那边去了？”}
}
{
    \eA{ask her what she's ailing from and what medicines she's taking, and explain
        to her that I know I ought to have gone over myself,}
}
{
}

\Verse{136}
{
    \zA{周瑞家的因说：“太太在那里，我回话去了， 姨太太就顺便叫我带来的。”}
}
{
    \eA{but that on my coming back from school a short while back, I again got a
        slight chill; and that I'll go in person another day." While Pao-yü was yet
        speaking, Hsi Hsüeh volunteered to take the message,}
}
{
}

\Verse{137}
{
    \zA{宝玉道：“宝姐姐在家里作什么呢？}
}
{
    \eA{and went off at once; and Mrs. Chou herself took her leave without another
        word. Mrs. Chou's son-in-law was,}
}
{
}

\Verse{138}
{
    \zA{怎么这几日也不 过来？”}
}
{
    \eA{in fact, Leng Tzu-hsing, the intimate friend of Yü-ts'un.}
}
{
}

\Verse{139}
{
    \zA{周瑞家的道：“身上不大好呢。”}
}
{
    \eA{Having recently become involved with some party in a lawsuit, on account of
        the sale of some curios,}
}
{
}

\Verse{140}
{
    \zA{宝玉听了，便和丫头们说：“谁去瞧瞧， 就说我和林姑娘打发来问姨娘姐姐安，问姐姐是什么病，吃什么药。}
}
{
    \eA{he had expressly charged his wife to come and sue for the favour (of a
        helping hand). Chou Jui's wife, relying upon her master's prestige, did not
        so much as take the affair to heart; and having waited till evening, she
        simply went over and requested lady Feng to befriend her,}
}
{
}

\Verse{141}
{
    \zA{论理，我该亲 自来的，就说才从学里回来，也着了些凉，改日再亲自来看。”}
}
{
    \eA{and the matter was forthwith ended. When the lamps were lit, lady Feng came
        over, after having disrobed herself, to see madame Wang. "I've already taken
        charge,"}
}
{
}

\Verse{142}
{
    \zA{说着，茜雪便答应 去了。}
}
{
    \eA{she observed, "of the things sent round to-day by the Chen family. As for
        the presents from us to them,}
}
{
}

\Verse{143}
{
    \zA{周瑞家的自去无话。}
}
{
    \eA{we should avail ourselves of the return of the boats,}
}
{
}

\Verse{144}
{
    \zA{原来周瑞家的女婿便是雨村的好友冷子兴，近日因卖古董，和人打官司，故叫 女人来讨情。}
}
{
    \eA{by which the fresh delicacies for the new year were forwarded, to hand them
        to them to carry back." Madame Wang nodded her head in token of approval.
        "The birthday presents," continued lady Feng, "for lady Ling Ngan,}
}
{
}

\Verse{145}
{
    \zA{周瑞家的仗着主子的势，把这些事也不放在心上，晚上只求求凤姐便 完了。}
}
{
    \eA{the mother of the Earl of Ling Ngan, have already been got together, and
        whom will you depute to take them over?" "See," suggested madame Wang, "who
        has nothing to do; let four maids go and all will be right!}
}
{
}

\Verse{146}
{
    \zA{至掌灯时，凤姐卸了妆，来见王夫人，回说：“今儿甄家送了来的东西，我已 收了。}
}
{
    \eA{why come again and ask me?" "Our eldest sister-in-law Chen," proceeded lady
        Feng, "came over to invite me to go to-morrow to their place for a little
        change. I don't think there will be anything for me to do to-morrow."}
}
{
}

\Verse{147}
{
    \zA{咱们送他的，趁着他家有年下送鲜的船，交给他带了去了。”}
}
{
    \eA{"Whether there be or not," replied madame Wang, "it doesn't matter; you must
        go, for whenever she comes with an invitation, it includes us, who are your
        seniors,}
}
{
}

\Verse{148}
{
    \zA{王夫人点点头 儿。}
}
{
    \eA{so that, of course, it isn't such a pleasant thing for you;}
}
{
}

\Verse{149}
{
    \zA{凤姐又道：“临安伯老太太生日的礼已经打点了，太太派谁送去？”}
}
{
    \eA{but as she doesn't ask us this time, but only asks you, it's evident that
        she's anxious that you should have a little distraction, and you mustn't
        disappoint her good intention.}
}
{
}

\Verse{150}
{
    \zA{王夫人道： “你瞧谁闲着，叫四个女人去就完了，又来问我。”}
}
{
    \eA{Besides it's certainly right that you should go over for a change." Lady
        Feng assented, and presently Li Wan, Ying Ch'un and the other cousins,}
}
{
}

\Verse{151}
{
    \zA{凤姐道：“今日珍大嫂子来请 我明日去逛逛，明日有什么事没有？”}
}
{
    \eA{likewise paid each her evening salutation and retired to their respective
        rooms, where nothing of any notice transpired. The next day lady Feng
        completed her toilette,}
}
{
}

\Verse{152}
{
    \zA{王夫人道：“有事没事都碍不着什么。}
}
{
    \eA{and came over first to tell madame Wang that she was off, and then went to
        say good-bye to dowager lady Chia;}
}
{
}

\Verse{153}
{
    \zA{每常 他来请，有我们，你自然不便；他不请我们单请你，可知是他的诚心叫你散荡散荡。}
}
{
    \eA{but when Pao-yü heard where she was going, he also wished to go; and as lady
        Feng had no help but to give in, and to wait until he had changed his
        clothes, the sister and brother-in-law got into a carriage,}
}
{
}

\Verse{154}
{
    \zA{别辜负了他的心，倒该过去走走才是。”}
}
{
    \eA{and in a short while entered the Ning mansion. Mrs. Yu, the wife of Chia
        Chen, and Mrs. Ch'in, the wife of Mr. Chia Jung,}
}
{
}

\Verse{155}
{
    \zA{凤姐答应了。}
}
{
    \eA{the two sisters-in-law, had,}
}
{
}

\Verse{156}
{
    \zA{当下李纨探春等姊妹们也都 定省毕，各归房无话。}
}
{
    \eA{along with a number of maids, waiting-girls, and other servants, come as far
        as the ceremonial gate to receive them,}
}
{
}

\Verse{157}
{
    \zA{次日凤姐梳洗了，先回王夫人毕，方来辞贾母。}
}
{
    \eA{and Mrs. Yu, upon meeting lady Feng, for a while indulged, as was her wont,
        in humorous remarks, after which,}
}
{
}

\Verse{158}
{
    \zA{宝玉听了，也要逛去，凤姐只 得答应着。}
}
{
    \eA{leading Pao-yü by the hand, they entered the drawing room and took their
        seats, Mrs. Ch'in handed tea round.}
}
{
}

\Verse{159}
{
    \zA{立等换了衣裳，姐儿两个坐了车。}
}
{
    \eA{"What have you people invited me to come here for?" promptly asked lady
        Feng; "if you have anything to present me with,}
}
{
}

\Verse{160}
{
    \zA{一时进入宁府，早有贾珍之妻尤氏与 贾蓉媳妇秦氏，婆媳两个带着多少侍妾丫鬟等接出仪门。}
}
{
    \eA{hand it to me at once, for I've other things to attend to." Mrs. Yu and Mrs.
        Ch'in had barely any time to exchange any further remarks, when several
        matrons interposed, smilingly: "Had our lady not come to-day, there would
        have been no help for it,}
}
{
}
\ChapterImage{img/chapters/023第7回 宴寧榮寶玉會秦鍾.jpg}


\Verse{161}
{
    \zA{那尤氏一见凤姐，必先嘲 笑一阵，一手拉了宝玉，同入上房里坐下。}
}
{
    \eA{but having come, you can't have it all your own way." While they were
        conversing about one thing and another, they caught sight of Chia Jung come
        in to pay his respects,}
}
{
}

\Verse{162}
{
    \zA{秦氏献了茶。}
}
{
    \eA{which prompted Pao-yü to inquire,}
}
{
}

\Verse{163}
{
    \zA{凤姐便说：“你们请我来 作什么？}
}
{
    \eA{"Isn't my elder brother at home to-day?" "He's gone out of town to-day,"}
}
{
}

\Verse{164}
{
    \zA{拿什么孝敬我？}
}
{
    \eA{replied Mrs. Yu, "to inquire after his grandfather.}
}
{
}

\Verse{165}
{
    \zA{有东西就献上来罢，我还有事呢！”}
}
{
    \eA{You'll find sitting here," she continued, "very dull, and why not go out and
        have a stroll?"}
}
{
}

\Verse{166}
{
    \zA{尤氏未及答应，几个媳 妇们先笑道：“二奶奶今日不来就罢，既来了，就依不得你老人家了。”}
}
{
    \eA{"A strange coincidence has taken place to-day," urged Mrs. Ch'in, with a
        smile; "some time back you, uncle Pao, expressed a wish to see my brother,}
}
{
}

\Verse{167}
{
    \zA{正说着， 只见贾蓉进来请安。}
}
{
    \eA{and to-day he too happens to be here at home. I think he's in the library;}
}
{
}

\Verse{168}
{
    \zA{宝玉因道：“大哥哥今儿不在家么？”}
}
{
    \eA{but why not go and see for yourself, uncle Pao?" Pao-yü descended at once
        from the stove-couch,}
}
{
}

\Verse{169}
{
    \zA{尤氏道：“今儿出城请 老爷的安去了。”}
}
{
    \eA{and was about to go, when Mrs. Yu bade the servants to mind and go with him.
        "Don't you let him get into trouble,"}
}
{
}

\Verse{170}
{
    \zA{又道：“可是你怪闷的，坐在这里作什么？}
}
{
    \eA{she enjoined. "It's a far different thing when he comes over under the
        charge of his grandmother,}
}
{
}

\Verse{171}
{
    \zA{何不出去逛逛呢？”}
}
{
    \eA{when he's all right." "If that be so,"}
}
{
}

\Verse{172}
{
    \zA{秦氏笑道：“今日可巧：上回宝二叔要见我兄弟，今儿他在这里书房里坐着呢，为 什么不瞧瞧去？”}
}
{
    \eA{remarked lady Feng, "why not ask the young gentleman to come in, and then I
        too can see him. There isn't, I hope, any objection to my seeing him?"
        "Never mind! never mind!" observed Mrs. Yu, smilingly; "it's as well that
        you shouldn't see him.}
}
{
}

\Verse{173}
{
    \zA{宝玉便去要见，尤氏忙吩咐人小心伺候着跟了去。}
}
{
    \eA{This brother of mine is not, like the boys of our Chia family, accustomed to
        roughly banging and knocking about.}
}
{
}

\Verse{174}
{
    \zA{凤姐道：“既 这么着，为什么不请进来我也见见呢？”}
}
{
    \eA{Other people's children are brought up politely and properly, and not in
        this vixenish style of yours. Why, you'd ridicule him to death!"}
}
{
}

\Verse{175}
{
    \zA{尤氏笑道：“罢，罢，可以不必见。}
}
{
    \eA{"I won't laugh at him then, that's all," smiled lady Feng; "tell them to
        bring him in at once."}
}
{
}

\Verse{176}
{
    \zA{比不 得咱们家的孩子，胡打海摔的惯了的。}
}
{
    \eA{"He's shy," proceeded Mrs. Ch'in, "and has seen nothing much of the world,
        so that you are sure to be put out when you see him,}
}
{
}

\Verse{177}
{
    \zA{人家的孩子都是斯斯文文的，没见过你这样 泼辣货。}
}
{
    \eA{sister." "What an idea!" exclaimed lady Feng. "Were he even No Cha himself,
        I'd like to see him; so don't talk trash; if, after all, you don't bring him
        round at once,}
}
{
}

\Verse{178}
{
    \zA{还叫人家笑话死呢！”}
}
{
    \eA{I'll give you a good slap on the mouth." "I daren't be obstinate,"}
}
{
}

\Verse{179}
{
    \zA{凤姐笑道：“我不笑话他就罢了，他敢笑话我？”}
}
{
    \eA{answered Mrs. Ch'in smiling; "I'll bring him round!" In a short while she
        did in fact lead in a young lad, who, compared with Pao-yü,}
}
{
}

\Verse{180}
{
    \zA{贾蓉道：“他生的腼腆，没见过大阵仗儿，婶子见了，没的生气。”}
}
{
    \eA{was somewhat more slight but, from all appearances, superior to Pao-yü in
        eyes and eyebrows, (good looks), which were so clear and well-defined, in
        white complexion and in ruddy lips,}
}
{
}

\Verse{181}
{
    \zA{凤姐啐道：“呸! 扯臊!他是哪吒我也要见见。}
}
{
    \eA{as well as graceful appearance and pleasing manners. He was however bashful
        and timid, like a girl. In a shy and demure way,}
}
{
}

\Verse{182}
{
    \zA{别放你娘的屁了!再不带来，打你顿好嘴巴子。”}
}
{
    \eA{he made a bow to lady Feng and asked after her health. Lady Feng was simply
        delighted with him. "You take a low seat next to him!"}
}
{
}

\Verse{183}
{
    \zA{贾蓉 溜湫着眼儿笑道：“何苦婶子又使利害!我们带了来就是了。”}
}
{
    \eA{she ventured laughingly as she first pushed Pao-yü back. Then readily
        stooping forward, she took this lad by the hand and asked him to take a seat
        next to her. Presently she inquired about his age,}
}
{
}

\Verse{184}
{
    \zA{凤姐也笑了。}
}
{
    \eA{his studies and such matters,}
}
{
}

\Verse{185}
{
    \zA{说着出去一会儿，果然带了个后生来：比宝玉略瘦些，眉清目秀，粉面朱唇， 身材俊俏，举止风流，似更在宝玉之上，只是怯怯羞羞有些女儿之态，腼腆含糊的
        向凤姐请安问好。}
}
{
    \eA{when she found that at school he went under the name of Ch'in Chung. The
        matrons and maids in attendance on lady Feng, perceiving that this was the
        first time their mistress met Ch'in Chung, (and knowing) that she had not at
        hand the usual presents, forthwith ran over to the other side and told P'ing
        Erh about it. P'ing Erh, aware of the close intimacy that existed between
        lady Feng and Mrs. Ch'in,}
}
{
}

\Verse{186}
{
    \zA{凤姐喜的先推宝玉笑道：“比下去了！”}
}
{
    \eA{speedily took upon herself to decide, and selecting a piece of silk, and two
        small gold medals,}
}
{
}

\Verse{187}
{
    \zA{便探身一把攥了这孩子 的手，叫他身旁坐下，慢慢问他年纪读书等事，方知他学名叫秦钟。}
}
{
    \eA{(bearing the wish that he should attain) the highest degree, the senior
        wranglership, she handed them to the servants who had come over, to take
        away. Lady Feng, however, explained that her presents were too mean by far,}
}
{
}

\Verse{188}
{
    \zA{早有凤姐跟的 丫鬟媳妇们，看见凤姐初见秦钟并未备得表礼来，遂忙过那边去告诉平儿。}
}
{
    \eA{but Mrs. Ch'in and the others expressed their appreciation of them; and in a
        short time the repast was over, and Mrs. Yu, lady Feng and Mrs. Ch'in played
        at dominoes, but of this no details need be given; while both Pao-yü and
        Ch'in Chung sat down,}
}
{
}

\Verse{189}
{
    \zA{平儿素 知凤姐和秦氏厚密，遂自作主意，拿了一匹尺头，两个“状元及第”的小金锞子， 交付来人送过去。}
}
{
    \eA{got up and talked, as they pleased. Since he had first glanced at Ch'in
        Chung, and seen what kind of person he was, he felt at heart as if he had
        lost something, and after being plunged in a dazed state for a time, he
        began again to give way to foolish thoughts in his mind.}
}
{
}

\Verse{190}
{
    \zA{凤姐还说太简薄些。}
}
{
    \eA{"There are then such beings as he in the world!"}
}
{
}

\Verse{191}
{
    \zA{秦氏等谢毕，一时吃过了饭，尤氏、凤姐、 秦氏等抹骨牌，不在话下。}
}
{
    \eA{he reflected. "I now see there are! I'm however no better than a wallowing
        pig or a mangy cow! Despicable destiny! why was I ever born in this
        household of a marquis and in the mansion of a duke?}
}
{
}

\Verse{192}
{
    \zA{宝玉、秦钟二人随便起坐说话儿。}
}
{
    \eA{Had I seen the light in the home of some penniless scholar, or
        poverty-stricken official,}
}
{
}

\Verse{193}
{
    \zA{那宝玉自一见秦钟，心中便如有所失，痴了 半日，自己心中又起了个呆想，乃自思道：“天下竟有这等的人物!如今看了，我
        竟成了泥猪癞狗了，可恨我为什么生在这侯门公府之家？}
}
{
    \eA{I could long ago have enjoyed the communion of his friendship, and I would
        not have lived my whole existence in vain! Though more honourable than he,
        it is indeed evident that silk and satins only serve to swathe this rotten
        trunk of mine, and choice wines and rich meats only to gorge the filthy
        drain and miry sewer of this body of mine! Wealth! and splendour! ye are no
        more than contaminated with pollution by me!"}
}
{
}

\Verse{194}
{
    \zA{要也生在寒儒薄宦的家 里，早得和他交接，也不枉生了一世。}
}
{
    \eA{Ever since Ch'in Chung had noticed Pao-yü's unusual appearance, his sedate
        deportment, and what is more, his hat ornamented with gold,}
}
{
}

\Verse{195}
{
    \zA{我虽比他尊贵，但绫锦纱罗，也不过裹了我 这枯株朽木；羊羔美酒，也不过填了我这粪窟泥沟。}
}
{
    \eA{and his dress full of embroidery, attended by beautiful maids and handsome
        youths, he did not indeed think it a matter of surprise that every one was
        fond of him. "Born as I have had the misfortune to be," he went on to
        commune within himself,}
}
{
}

\Verse{196}
{
    \zA{‘富贵’二字，真真把人荼毒 了。”}
}
{
    \eA{"in an honest, though poor family, how can I presume to enjoy his
        companionship!}
}
{
}

\Verse{197}
{
    \zA{那秦钟见了宝玉形容出众，举止不凡，更兼金冠绣服，艳婢娇童，――“果 然怨不得姐姐素日提起来就夸不绝口。}
}
{
    \eA{This is verily a proof of what a barrier poverty and wealth set between man
        and man. What a serious misfortune is this too in this mortal world!" In
        wild and inane ideas of the same strain, indulged these two youths! Pao-yü
        by and by further asked of him what books he was reading,}
}
{
}

\Verse{198}
{
    \zA{我偏偏生于清寒之家，怎能和他交接亲厚一 番，也是缘法”。}
}
{
    \eA{and Ch'in Chung, in answer to these inquiries, told him the truth. A few
        more questions and answers followed; and after about ten remarks,}
}
{
}

\Verse{199}
{
    \zA{二人一样胡思乱想。}
}
{
    \eA{a greater intimacy sprang up between them.}
}
{
}

\Verse{200}
{
    \zA{宝玉又问他读什么书，秦钟见问，便依实而 答。}
}
{
    \eA{Tea and fruits were shortly served, and while they were having their tea,
        Pao-yü suggested, "We two don't take any wine,}
}
{
}

\Verse{201}
{
    \zA{二人你言我语，十来句话，越觉亲密起来了。}
}
{
    \eA{and why shouldn't we have our fruit served on the small couch inside, and go
        and sit there, and thus save you all the trouble?"}
}
{
}

\Verse{202}
{
    \zA{一时捧上茶果吃茶，宝玉便说： “我们两个又不吃酒，把果子摆在里间小炕上，我们那里去，省了闹的你们不安。”}
}
{
    \eA{The two of them thereupon came into the inner apartment to have their tea;
        and Mrs. Ch'in attended to the laying out of fruit and wines for lady Feng,
        and hurriedly entered the room and hinted to Pao-yü: "Dear uncle Pao, your
        nephew is young, and should he happen to say anything disrespectful,}
}
{
}

\Verse{203}
{
    \zA{于是二人进里间来吃茶。}
}
{
    \eA{do please overlook it, for my sake, for though shy, he's naturally of a
        perverse and wilful disposition,}
}
{
}

\Verse{204}
{
    \zA{秦氏一面张罗凤姐吃果酒，一面忙进来嘱咐宝玉道：“宝 二叔：你侄儿年轻，倘或说话不防头，你千万看着我，别理他。}
}
{
    \eA{and is rather given to having his own way." "Off with you!" cried Pao-yü
        laughing; "I know it all." Mrs. Ch'in then went on to give a bit of advice
        to her brother, and at length came to keep lady Feng company. Presently lady
        Feng and Mrs. Yu sent another servant to tell Pao-yü that there was outside
        of everything they might wish to eat and that they should mind and go and
        ask for it;}
}
{
}

\Verse{205}
{
    \zA{他虽腼腆，却脾气 拐孤，不大随和儿。”}
}
{
    \eA{and Pao-yü simply signified that they would; but his mind was not set upon
        drinking or eating;}
}
{
}

\Verse{206}
{
    \zA{宝玉笑道：“你去罢，我知道了。”}
}
{
    \eA{all he did was to keep making inquiries of Ch'in Chung about recent family
        concerns.}
}
{
}

\Verse{207}
{
    \zA{秦氏又嘱咐了他兄弟一 回，方去陪凤姐儿去了。}
}
{
    \eA{Ch'in Chung went on to explain that his tutor had last year relinquished his
        post, that his father was advanced in years and afflicted with disease,}
}
{
}

\Verse{208}
{
    \zA{一时凤姐尤氏又打发人来问宝玉：“要吃什么，只管要去。”}
}
{
    \eA{and had multifarious public duties to preoccupy his mind, so that he had as
        yet had no time to make arrangements for another tutor, and that all he did
        was no more than to keep up his old tasks;}
}
{
}

\Verse{209}
{
    \zA{宝玉只答应着， 也无心在饮食上，只问秦钟近日家务等事。}
}
{
    \eA{that as regards study, it was likewise necessary to have the company of one
        or two intimate friends, as then only, by dint of a frequent exchange of
        ideas and opinions,}
}
{
}

\Verse{210}
{
    \zA{秦钟因言：“业师于去岁辞馆，家父年 纪老了，残疾在身，公务繁冗，因此尚未议及延师，目下不过在家温习旧课而已。}
}
{
    \eA{one could arrive at progress; and Pao-yü gave him no time to complete, but
        eagerly urged, "Quite so! But in our household, we have a family school, and
        those of our kindred who have no means sufficient to engage the services of
        a tutor are at liberty to come over for the sake of study,}
}
{
}

\Verse{211}
{
    \zA{再读书一事也必须有一二知己为伴，时常大家讨论才能有些进益――”宝玉不待说 完，便道：“正是呢!我们家却有个家塾，合族中有不能延师的便可入塾读书，亲
        戚子弟可以附读。}
}
{
    \eA{and the sons and brothers of our relatives are likewise free to join the
        class. As my own tutor went home last year, I am now also wasting my time
        doing nothing; my father's intention was that I too should have gone over to
        this school, so that I might at least temporarily keep up what I have
        already read, pending the arrival of my tutor next year,}
}
{
}

\Verse{212}
{
    \zA{我因上年业师回家去了，也现荒废着。}
}
{
    \eA{when I could again very well resume my studies alone at home. But my
        grandmother raised objections;}
}
{
}

\Verse{213}
{
    \zA{家父之意亦欲暂送我去， 且温习着旧书，待明年业师上来，再各自在家读书。}
}
{
    \eA{maintaining first of all, that the boys who attend the family classes being
        so numerous, she feared we would be sure to be up to mischief, which
        wouldn't be at all proper; and that, in the second place,}
}
{
}

\Verse{214}
{
    \zA{家祖母因说：一则家学里子弟 太多，恐怕大家淘气，反不好；二则也因我病了几天，遂暂且耽搁着。}
}
{
    \eA{as I had been ill for some time, the matter should be dropped, for the
        present. But as, from what you say, your worthy father is very much
        exercised on this score, you should, on your return, tell him all about it,
        and come over to our school.}
}
{
}
\ChapterImage{img/chapters/024第7回 鳳姐坐車聞焦大罵.jpg}


\Verse{215}
{
    \zA{如此说来， 尊翁如今也为此事悬心，今日回去，何不禀明，就在我们这敝塾中来？}
}
{
    \eA{I'll also be there as your schoolmate; and as you and I will reap mutual
        benefit from each other's companionship, won't it be nice!" "When my father
        was at home the other day," Ch'in Chung smiled and said,}
}
{
}

\Verse{216}
{
    \zA{我也相伴， 彼此有益，岂不是好事？”}
}
{
    \eA{"he alluded to the question of a tutor, and explained that the free schools
        were an excellent institution.}
}
{
}

\Verse{217}
{
    \zA{秦钟笑道：“家父前日在家提起延师一事，也曾提起这 里的义学倒好，原要来和这里的老爷商议引荐；因这里又有事忙，不便为这点子小 事来絮聒。}
}
{
    \eA{He even meant to have come and talked matters over with his son-in-law's
        father about my introduction, but with the urgent concerns here, he didn't
        think it right for him to come about this small thing, and make any trouble.
        But if you really believe that I might be of use to you, in either grinding
        the ink, or washing the slab, why shouldn't you at once make the needful
        arrangements,}
}
{
}

\Verse{218}
{
    \zA{二叔果然度量侄儿或可磨墨洗砚，何不速速作成，彼此不致荒废，既可 以常相聚谈，又可以慰父母之心，又可以得朋友之乐，岂不是美事？”}
}
{
    \eA{so that neither you nor I may idle our time? And as we shall be able to come
        together often and talk matters over, and set at the same time our parents'
        minds at ease, and to enjoy the pleasure of friendship, won't it be a
        profitable thing!" "Compose your mind!" suggested Pao-yü.}
}
{
}

\Verse{219}
{
    \zA{宝玉道：“放 心，放心!咱们回来告诉你姐夫姐姐和琏二嫂子，今日你就回家禀明令尊，我回去 禀明了祖母，再无不速成之理。”}
}
{
    \eA{"We can by and by first of all, tell your brother-in-law, and your sister as
        well as sister-in-law Secunda Lien; and on your return home to-day, lose no
        time in explaining all to your worthy father, and when I get back, I'll
        speak to my grandmother; and I can't see why our wishes shouldn't speedily
        be accomplished."}
}
{
}

\Verse{220}
{
    \zA{二人计议已定，那天气已是掌灯时分，出来又看他们玩了一回牌。}
}
{
    \eA{By the time they had arrived at this conclusion, the day was far advanced,
        and the lights were about to be lit; and they came out and watched them once
        more for a time as they played at dominoes.}
}
{
}

\Verse{221}
{
    \zA{算帐时，却 又是秦氏尤氏二人输了戏酒的东道，言定后日吃这东道，一面又吃了晚饭。}
}
{
    \eA{When they came to settle their accounts Mrs. Ch'in and Mrs. Yu were again
        the losers and had to bear the expense of a theatrical and dinner party; and
        while deciding that they should enjoy this treat the day after the morrow,
        they also had the evening repast.}
}
{
}

\Verse{222}
{
    \zA{因天黑 了，尤氏说：“派两个小子送了秦哥儿家去。”}
}
{
    \eA{Darkness having set in, Mrs. Yu gave orders that two youths should accompany
        Mr. Ch'in home. The matrons went out to deliver the directions, and after a
        somewhat long interval,}
}
{
}

\Verse{223}
{
    \zA{媳妇们传出去半日。}
}
{
    \eA{Ch'in Chung said goodbye and was about to start on his way.}
}
{
}

\Verse{224}
{
    \zA{秦钟告辞起身， 尤氏问：“派谁送去？”}
}
{
    \eA{"Whom have you told off to escort him?" asked Mrs. Yu. "Chiao Ta," replied
        the matrons, "has been told to go,}
}
{
}

\Verse{225}
{
    \zA{媳妇们回说：“外头派了焦大，谁知焦大醉了，又骂呢。”}
}
{
    \eA{but it happens that he's under the effects of drink and making free use
        again of abusive language." Mrs. Yu and Mrs. Chin remonstrated.}
}
{
}

\Verse{226}
{
    \zA{尤氏秦氏都道：“偏又派他作什么？}
}
{
    \eA{"What's the use," they said, "of asking him? that mean fellow shouldn't be
        chosen,}
}
{
}

\Verse{227}
{
    \zA{那个小子派不得？}
}
{
    \eA{but you will go again and provoke him."}
}
{
}

\Verse{228}
{
    \zA{偏又惹他！”}
}
{
    \eA{"People always maintain,"}
}
{
}

\Verse{229}
{
    \zA{凤姐道：“成日 家说你太软弱了，纵的家里人这样，还了得吗？”}
}
{
    \eA{added lady Feng, "that you are far too lenient. But fancy allowing servants
        in this household to go on in this way; why, what will be the end of it?"}
}
{
}

\Verse{230}
{
    \zA{尤氏道：“你难道不知这焦大的？}
}
{
    \eA{"You don't mean to tell me," observed Mrs. Yu, "that you don't know this
        Chiao Ta?}
}
{
}

\Verse{231}
{
    \zA{连老爷都不理他，你珍大哥哥也不理他。}
}
{
    \eA{Why, even the gentlemen one and all pay no heed to his doings! your eldest
        brother, Chia Cheng, he too doesn't notice him.}
}
{
}

\Verse{232}
{
    \zA{因他从小儿跟着太爷出过三四回兵，从死 人堆里把太爷背出来了，才得了命；自己挨着饿，却偷了东西给主子吃；两日没水，
        得了半碗水，给主子喝，他自己喝马溺：不过仗着这些功劳情分，有祖宗时，都另 眼相待，如今谁肯难为他？}
}
{
    \eA{It's all because when he was young he followed our ancestor in three or four
        wars, and because on one occasion, by extracting our senior from the heap of
        slain and carrying him on his back, he saved his life. He himself suffered
        hunger and stole food for his master to eat; they had no water for two days;}
}
{
}

\Verse{233}
{
    \zA{他自己又老了，又不顾体面，一味的好酒，喝醉了无人 不骂。}
}
{
    \eA{and when he did get half a bowl, he gave it to his master, while he himself
        had sewage water. He now simply presumes upon the sentimental obligations
        imposed by these services.}
}
{
}

\Verse{234}
{
    \zA{我常说给管事的，以后不用派他差使，只当他是个死的就完了。}
}
{
    \eA{When the seniors of the family still lived, they all looked upon him with
        exceptional regard; but who at present ventures to interfere with him? He is
        also advanced in years,}
}
{
}

\Verse{235}
{
    \zA{今儿又派了 他！”}
}
{
    \eA{and doesn't care about any decent manners;}
}
{
}

\Verse{236}
{
    \zA{凤姐道：“我何曾不知这焦大？}
}
{
    \eA{his sole delight is wine; and when he gets drunk, there isn't a single
        person whom he won't abuse.}
}
{
}

\Verse{237}
{
    \zA{到底是你们没主意，何不远远的打发他到庄 子上去就完了！”}
}
{
    \eA{I've again and again told the stewards not to henceforward ask Chiao Ta to
        do any work whatever, but to treat him as dead and gone; and here he's sent
        again to-day." "How can I not know all about this Chiao Ta?"}
}
{
}

\Verse{238}
{
    \zA{说着，因问：“我们的车可齐备了？”}
}
{
    \eA{remarked lady Feng; "but the secret of all this trouble is, that you won't
        take any decisive step.}
}
{
}

\Verse{239}
{
    \zA{众媳妇们说：“伺候齐了。”}
}
{
    \eA{Why not pack him off to some distant farm, and have done with him?"}
}
{
}

\Verse{240}
{
    \zA{凤姐也起身告辞，和宝玉携手同行。}
}
{
    \eA{And as she spoke, "Is our carriage ready?" she went on to inquire. "All
        ready and waiting,"}
}
{
}

\Verse{241}
{
    \zA{尤氏等送至大厅前，见灯火辉煌，众小厮 都在丹墀侍立。}
}
{
    \eA{interposed the married women. Lady Feng also got up, said good-bye, and hand
        in hand with Pao-yü, they walked out of the room, escorted by Mrs. Yu and
        the party,}
}
{
}

\Verse{242}
{
    \zA{那焦大又恃贾珍不在家，因趁着酒兴，先骂大总管赖二，说他：“不 公道，欺软怕硬!有好差使派了别人，这样黑更半夜送人就派我，没良心的忘八羔
        子!瞎充管家!你也不想想焦大太爷跷起一只腿，比你的头还高些。}
}
{
    \eA{as far as the entrance of the Main Hall, where they saw the lamps shedding a
        brilliant light and the attendants all waiting on the platforms. Chiao Ta,
        however, availing himself of Chia Chen's absence from home, and elated by
        wine, began to abuse the head steward Lai Erh for his injustice. "You bully
        of the weak and coward with the strong," he cried, "when there's any
        pleasant charge,}
}
{
}

\Verse{243}
{
    \zA{二十年头里的焦 大太爷眼里有谁？}
}
{
    \eA{you send the other servants, but when it's a question of seeing any one home
        in the dark, then you ask me, you disorderly clown!}
}
{
}

\Verse{244}
{
    \zA{别说你们这一把子的杂种们！”}
}
{
    \eA{a nice way you act the steward, indeed! Do you forget that if Mr. Chiao Ta
        chose to raise one leg,}
}
{
}

\Verse{245}
{
    \zA{正骂得兴头上，贾蓉送凤姐的车 出来。}
}
{
    \eA{it would be a good deal higher than your head! Remember please, that twenty
        years ago, Mr. Chiao Ta wouldn't even so much as look at any one,}
}
{
}

\Verse{246}
{
    \zA{众人喝他不住，贾蓉忍不住便骂了几句，叫人：“捆起来!等明日酒醒了， 再问他还寻死不寻死！”}
}
{
    \eA{no matter who it was; not to mention a pack of hybrid creatures like
        yourselves!" While he went on cursing and railing with all his might, Chia
        Jung appeared walking by lady Feng's carriage. All the servants having tried
        to hush him and not succeeding,}
}
{
}

\Verse{247}
{
    \zA{那焦大那里有贾蓉在眼里？}
}
{
    \eA{Chia Jung became exasperated; and forthwith blew him up for a time.}
}
{
}

\Verse{248}
{
    \zA{反大叫起来，赶着贾蓉叫：“蓉 哥儿，你别在焦大跟前使主子性儿!别说你这样儿的，就是你爹、你爷爷，也不敢 和焦大挺腰子呢。}
}
{
    \eA{"Let some one bind him up," he cried, "and tomorrow, when he's over the
        wine, I'll call him to task, and we'll see if he won't seek death." Chiao Ta
        showed no consideration for Chia Jung.}
}
{
}

\Verse{249}
{
    \zA{不是焦大一个人，你们作官儿，享荣华，受富贵!你祖宗九死一 生挣下这个家业，到如今不报我的恩，反和我充起主子来了。}
}
{
    \eA{On the contrary, he shouted with more vigour. Going up to Chia Jung:
        "Brother Jung," he said, "don't put on the airs of a master with Chiao Ta.
        Not to speak of a man such as you, why even your father and grandfather
        wouldn't presume to display such side with Chiao Ta.}
}
{
}

\Verse{250}
{
    \zA{不和我说别的还可； 再说别的，咱们‘白刀子进去，红刀子出来’！”}
}
{
    \eA{Were it not for Chiao Ta, and him alone, where would your office, honours,
        riches and dignity be? Your ancestor, whom I brought back from the jaws of
        death,}
}
{
}

\Verse{251}
{
    \zA{凤姐在车上和贾蓉说：“还不早 些打发了没王法的东西!留在家里，岂不是害？}
}
{
    \eA{heaped up all this estate, but up to this very day have I received no thanks
        for the services I rendered! on the contrary, you come here and play the
        master; don't say a word more, and things may come right;}
}
{
}

\Verse{252}
{
    \zA{亲友知道，岂不笑话咱们这样的人 家，连个规矩都没有？”}
}
{
    \eA{but if you do, I'll plunge the blade of a knife white in you and extract it
        red." Lady Feng, from inside the carriage, remarked to Chia Jung: "Don't you
        yet pack off this insolent fellow!}
}
{
}

\Verse{253}
{
    \zA{贾蓉答应了“是”。}
}
{
    \eA{Why, if you keep him in your house, won't he be a source of mischief?}
}
{
}

\Verse{254}
{
    \zA{众人见他太撒野，只得上来了几个，揪翻捆倒，拖往马圈里去。}
}
{
    \eA{Besides, were relatives and friends to hear about these things, won't they
        have a laugh at our expense, that a household like ours should be so devoid
        of all propriety?"}
}
{
}

\Verse{255}
{
    \zA{焦大益发连贾 珍都说出来，乱嚷乱叫，说：“要往祠堂里哭太爷去，那里承望到如今生下这些畜
        生来!每日偷狗戏鸡，爬灰的爬灰，养小叔子的养小叔子，我什么不知道？}
}
{
    \eA{Chia Jung assented. The whole band of servants finding that Chiao Ta was
        getting too insolent had no help but to come up and throw him over, and
        binding him up, they dragged him towards the stables. Chiao Ta abused even
        Chia Chen with still more vehemence, and shouted in a boisterous manner.}
}
{
}

\Verse{256}
{
    \zA{咱们‘胳 膊折了往袖子里藏’！”}
}
{
    \eA{"I want to go," he cried, "to the family Ancestral Temple and mourn my old
        master.}
}
{
}

\Verse{257}
{
    \zA{众小厮见说出来的话有天没日的，唬得魂飞魄丧，把他捆 起来，用土和马粪满满的填了他一嘴。}
}
{
    \eA{Who would have ever imagined that he would leave behind such vile creatures
        of descendants as you all, day after day indulging in obscene and incestuous
        practices, 'in scraping of the ashes' and in philandering with
        brothers-in-law. I know all about your doings; the best thing is to hide
        one's stump of an arm in one's sleeve!"}
}
{
}

\Verse{258}
{
    \zA{凤姐和贾蓉也遥遥的听见了，都装作没听见。}
}
{
    \eA{(wash one's dirty clothes at home). The servants who stood by, upon hearing
        this wild talk, were quite at their wits' end,}
}
{
}

\Verse{259}
{
    \zA{宝玉在车上听见，因问凤姐道： “姐姐，你听他说‘爬灰的爬灰’，这是什么话？”}
}
{
    \eA{and they at once seized him, tied him up, and filled his mouth to the
        fullest extent with mud mixed with some horse refuse. Lady Feng and Chia
        Jung heard all he said from a distance, but pretended not to hear;}
}
{
}

\Verse{260}
{
    \zA{凤姐连忙喝道：“少胡说!那 是醉汉嘴里胡？}
}
{
    \eA{but Pao-yü, seated in the carriage as he was, also caught this extravagant
        talk and inquired of lady Feng:}
}
{
}

\Verse{261}
{
    \zA{，你是什么样的人，不说没听见，还倒细问!等我回了太太，看是 捶你不捶你！”}
}
{
    \eA{"Sister, did you hear him say something about 'scraping of the ashes?'
        What's it?" "Don't talk such rubbish!" hastily shouted lady Feng; "it was
        the maudlin talk of a drunkard! A nice boy you are! not to speak of your
        listening,}
}
{
}

\Verse{262}
{
    \zA{吓得宝玉连忙央告：“好姐姐，我再不敢说这些话了。”}
}
{
    \eA{but you must also inquire! wait and I'll tell your mother and we'll see if
        she doesn't seriously take you to task." Pao-yü was in such a state of
        fright that he speedily entreated her to forgive him.}
}
{
}

\Verse{263}
{
    \zA{凤姐哄他 道：“好兄弟，这才是呢。}
}
{
    \eA{"My dear sister," he craved, "I won't venture again to say anything of the
        kind" "My dear brother,}
}
{
}

\Verse{264}
{
    \zA{等回去咱们回了老太太，打发人到家学里去说明了，请 了秦钟学里念书去要紧。”}
}
{
    \eA{if that be so, it's all right!" rejoined lady Feng reassuringly; "on our
        return we'll speak to her venerable ladyship and ask her to send some one to
        arrange matters in the family school, and invite Ch'in Chung to come to
        school for his studies."}
}
{
}

\Verse{265}
{
    \zA{说着自回荣府而来。}
}
{
    \eA{While yet this conversation was going on,}
}
{
}

\Verse{266}
{
    \zA{要知端的，下回分解。}
}
{
    \eA{they arrived at the Jung Mansion. Reader, do you wish to know what follows?}
}
{
}

\Verse{267}
{
    \zA{(本章完)}
}
{
    \eA{if you do, the next chapter will unfold it.}
}
{
}
